% Options for packages loaded elsewhere
\PassOptionsToPackage{unicode}{hyperref}
\PassOptionsToPackage{hyphens}{url}
\documentclass[
]{article}
\usepackage{xcolor}
\usepackage{amsmath,amssymb}
\setcounter{secnumdepth}{-\maxdimen} % remove section numbering
\usepackage{iftex}
\ifPDFTeX
  \usepackage[T1]{fontenc}
  \usepackage[utf8]{inputenc}
  \usepackage{textcomp} % provide euro and other symbols
\else % if luatex or xetex
  \usepackage{unicode-math} % this also loads fontspec
  \defaultfontfeatures{Scale=MatchLowercase}
  \defaultfontfeatures[\rmfamily]{Ligatures=TeX,Scale=1}
\fi
\usepackage{lmodern}
\ifPDFTeX\else
  % xetex/luatex font selection
\fi
% Use upquote if available, for straight quotes in verbatim environments
\IfFileExists{upquote.sty}{\usepackage{upquote}}{}
\IfFileExists{microtype.sty}{% use microtype if available
  \usepackage[]{microtype}
  \UseMicrotypeSet[protrusion]{basicmath} % disable protrusion for tt fonts
}{}
\makeatletter
\@ifundefined{KOMAClassName}{% if non-KOMA class
  \IfFileExists{parskip.sty}{%
    \usepackage{parskip}
  }{% else
    \setlength{\parindent}{0pt}
    \setlength{\parskip}{6pt plus 2pt minus 1pt}}
}{% if KOMA class
  \KOMAoptions{parskip=half}}
\makeatother
\setlength{\emergencystretch}{3em} % prevent overfull lines
\providecommand{\tightlist}{%
  \setlength{\itemsep}{0pt}\setlength{\parskip}{0pt}}
\usepackage{bookmark}
\IfFileExists{xurl.sty}{\usepackage{xurl}}{} % add URL line breaks if available
\urlstyle{same}
\hypersetup{
  pdftitle={Advanced AI-Powered Development},
  hidelinks,
  pdfcreator={LaTeX via pandoc}}

\title{Advanced \texttt{AI}-Powered Development}
\author{}
\date{}

\begin{document}
\maketitle

\subsection{Description}\label{description}

\texttt{AI}-powered development has rapidly evolved beyond simple code
completion into a rich ecosystem of agents, tools, skills, \texttt{MCP}
servers, \texttt{RAG} pipelines, and autonomous coding workflows. Modern
developers must understand how to harness these capabilities to work
effectively with large codebases, monoliths, and complex systems. This
course provides deep, hands-on coverage of the full \texttt{AI}
development stack as it exists today, from terminal-based \texttt{AI}
coding agents to building custom \texttt{MCP} servers and
\texttt{RAG}-augmented development workflows.

The course is designed for experienced developers who already use basic
\texttt{AI} assistance and want to move to the next level: working with
agentic tools, managing context in massive codebases, building custom
integrations, and understanding the architecture behind modern
\texttt{AI} development infrastructure.

All topics reflect the state of the art as of early 2026, covering the
latest tools, protocols, and patterns that are actively used in
production environments.

\subsection{Duration}\label{duration}

40 hours / 5 days

\subsection{Intended Audience}\label{intended-audience}

\begin{itemize}
\tightlist
\item
  Senior developers who want to master advanced \texttt{AI}-powered
  development workflows
\item
  Software architects evaluating \texttt{AI} tool integration strategies
  for large systems
\item
  Tech leads establishing \texttt{AI} development practices for their
  teams
\item
  \texttt{DevOps} engineers seeking \texttt{AI}-powered automation for
  complex infrastructure
\item
  Developers working with large codebases or monoliths who need
  effective \texttt{AI} strategies
\end{itemize}

\subsection{Prerequisites}\label{prerequisites}

\begin{itemize}
\tightlist
\item
  \texttt{Solid} experience in at least one modern programming language
  (\texttt{Python}, \texttt{JavaScript}/\texttt{TypeScript},
  \texttt{Java}, Go, or similar)
\item
  Working knowledge of \texttt{Git}, \texttt{CI/CD}, and modern
  development workflows
\item
  Basic experience using \texttt{AI} coding assistants (e.g.,
  \texttt{GitHub\ Copilot}, \texttt{ChatGPT}, Claude)
\item
  Familiarity with APIs, \texttt{REST}, and \texttt{JSON}
\item
  Understanding of software architecture concepts
  (\texttt{microservices}, monoliths, APIs)
\end{itemize}

\subsection{Required Knowledge}\label{required-knowledge}

\begin{itemize}
\tightlist
\item
  \texttt{Python} Programming (or equivalent experience)
\item
  \texttt{HTML} \texttt{CSS} \texttt{JavaScript} Fundamentals (or
  equivalent experience)
\item
  \texttt{TypeScript} Programming (or equivalent experience)
\end{itemize}

\subsection{Objectives}\label{objectives}

\begin{itemize}
\tightlist
\item
  Effectively use agentic \texttt{AI} coding tools to navigate,
  understand, and modify large codebases
\item
  Build and deploy custom \texttt{MCP} servers to extend \texttt{AI}
  assistants with project-specific capabilities
\item
  Design and implement \texttt{RAG} pipelines to give \texttt{AI} tools
  deep knowledge of proprietary codebases and documentation
\item
  Understand and apply the tool-use, skills, and agent patterns in
  modern \texttt{AI} development
\item
  Develop strategies for \texttt{AI}-assisted work on monoliths and
  large-scale systems
\item
  Evaluate and integrate the latest \texttt{AI} development tools and
  frameworks
\item
  Build custom agents for specialized development tasks
\end{itemize}

\subsection{Outline}\label{outline}

\begin{itemize}
\tightlist
\item
  The Modern \texttt{AI} Development Landscape

  \begin{itemize}
  \tightlist
  \item
    Evolution from code completion to agentic development
  \item
    Categories of \texttt{AI} development tools

    \begin{itemize}
    \tightlist
    \item
      Inline completion engines
    \item
      Chat-based assistants
    \item
      Terminal-based agents
    \item
      IDE-integrated agents
    \item
      Autonomous coding agents
    \end{itemize}
  \item
    Key concepts: context \texttt{windows}, token limits, and their
    practical implications
  \item
    The role of models: Claude, \texttt{GPT}, Gemini, open-source models
  \item
    Cloud-hosted vs local models for development
  \item
    Cost, latency, and quality tradeoffs
  \item
    Privacy and intellectual property considerations
  \end{itemize}
\item
  \texttt{AI} Coding Agents

  \begin{itemize}
  \tightlist
  \item
    What makes an agent different from a chatbot
  \item
    Agent architectures: plan-execute, \texttt{ReAct}, tool-use loops
  \item
    Popular agent-based tools

    \begin{itemize}
    \tightlist
    \item
      \texttt{Claude\ Code} (\texttt{CLI} agent)
    \item
      Cursor (agent mode)
    \item
      \texttt{Windsurf}
    \item
      \texttt{Aider}
    \item
      \texttt{Cline}
    \item
      \texttt{Copilot\ Workspace}
    \end{itemize}
  \item
    Agent capabilities

    \begin{itemize}
    \tightlist
    \item
      \texttt{File} reading and editing
    \item
      Running shell commands
    \item
      Searching codebases
    \item
      Running tests and interpreting results
    \item
      \texttt{Git} operations
    \item
      Multi-step task execution
    \end{itemize}
  \item
    Working with agent permissions and safety models
  \item
    Configuring agents for your project (.\texttt{claude}/,
    \texttt{CLAUDE}.md, rules files)
  \item
    Giving agents project context and conventions
  \item
    Comparing agent effectiveness across tasks
  \end{itemize}
\item
  Working with Large Codebases and Monoliths

  \begin{itemize}
  \tightlist
  \item
    Challenges of \texttt{AI} in large codebases

    \begin{itemize}
    \tightlist
    \item
      Context window limitations
    \item
      Navigating unfamiliar code
    \item
      Understanding implicit conventions
    \item
      Cross-module dependencies
    \end{itemize}
  \item
    Strategies for effective \texttt{AI} use in monoliths

    \begin{itemize}
    \tightlist
    \item
      Structured project documentation for \texttt{AI} consumption
    \item
      Writing effective \texttt{CLAUDE}.md / rules files
    \item
      Modular context feeding
    \item
      Incremental exploration patterns
    \end{itemize}
  \item
    Code navigation with \texttt{AI} agents

    \begin{itemize}
    \tightlist
    \item
      Searching by pattern and semantic meaning
    \item
      Understanding call graphs and dependencies
    \item
      Tracing data flow across modules
    \end{itemize}
  \item
    Refactoring monoliths with \texttt{AI} assistance

    \begin{itemize}
    \tightlist
    \item
      Identifying extraction candidates
    \item
      Safe, incremental refactoring workflows
    \item
      Maintaining test coverage during refactoring
    \end{itemize}
  \item
    \texttt{AI}-assisted code archaeology

    \begin{itemize}
    \tightlist
    \item
      Understanding legacy code
    \item
      Documenting tribal knowledge
    \item
      Mapping undocumented APIs and protocols
    \end{itemize}
  \end{itemize}
\item
  Tool Use and Function Calling

  \begin{itemize}
  \tightlist
  \item
    The tool-use pattern in \texttt{LLMs}

    \begin{itemize}
    \tightlist
    \item
      How \texttt{LLMs} invoke external tools
    \item
      Tool definitions and schemas
    \item
      Tool result handling
    \item
      Multi-turn tool-use conversations
    \end{itemize}
  \item
    Designing effective tools for \texttt{AI} agents

    \begin{itemize}
    \tightlist
    \item
      Tool granularity and composability
    \item
      Input/output schema design
    \item
      Error handling and feedback
    \item
      Idempotency and safety
    \end{itemize}
  \item
    Built-in tools in popular agents

    \begin{itemize}
    \tightlist
    \item
      \texttt{File} operations (read, write, edit, glob, \texttt{grep})
    \item
      Shell execution
    \item
      Web search and fetch
    \item
      LSP integration
    \end{itemize}
  \item
    Building custom tools

    \begin{itemize}
    \tightlist
    \item
      Tool definition formats
    \item
      Implementing tool handlers
    \item
      Testing tools with \texttt{AI} agents
    \item
      Versioning and maintaining tools
    \end{itemize}
  \end{itemize}
\item
  \texttt{MCP} (Model Context Protocol)

  \begin{itemize}
  \tightlist
  \item
    What is \texttt{MCP} and why it matters

    \begin{itemize}
    \tightlist
    \item
      The problem \texttt{MCP} solves
    \item
      \texttt{MCP} architecture: hosts, clients, servers
    \item
      Transport mechanisms (stdio, \texttt{SSE}, streamable
      \texttt{HTTP})
    \end{itemize}
  \item
    \texttt{MCP} capabilities

    \begin{itemize}
    \tightlist
    \item
      Tools: actions the \texttt{AI} can perform
    \item
      Resources: data the \texttt{AI} can read
    \item
      Prompts: reusable prompt templates
    \item
      Sampling: letting servers request \texttt{LLM} completions
    \end{itemize}
  \item
    Using existing \texttt{MCP} servers

    \begin{itemize}
    \tightlist
    \item
      Filesystem server
    \item
      \texttt{GitHub} server
    \item
      Database servers (\texttt{PostgreSQL}, \texttt{SQLite})
    \item
      Slack, \texttt{Jira}, Linear servers
    \item
      Browser automation servers (\texttt{Puppeteer},
      \texttt{Playwright})
    \end{itemize}
  \item
    Building custom \texttt{MCP} servers

    \begin{itemize}
    \tightlist
    \item
      Server implementation in \texttt{Python} and \texttt{TypeScript}
    \item
      Exposing project-specific tools and resources
    \item
      Authentication and authorization
    \item
      Deployment and lifecycle management
    \end{itemize}
  \item
    \texttt{MCP} in team environments

    \begin{itemize}
    \tightlist
    \item
      Shared \texttt{MCP} server configurations
    \item
      Central tool registries
    \item
      Security policies and access control
    \end{itemize}
  \end{itemize}
\item
  Skills and Slash Commands

  \begin{itemize}
  \tightlist
  \item
    What are skills in \texttt{AI} development tools
  \item
    Built-in skills vs custom skills
  \item
    Creating custom skills

    \begin{itemize}
    \tightlist
    \item
      Skill definition and structure
    \item
      Parameterized skills
    \item
      Skill composition and chaining
    \end{itemize}
  \item
    Practical skill examples

    \begin{itemize}
    \tightlist
    \item
      Code review skills
    \item
      Commit message generation
    \item
      PR creation workflows
    \item
      Deployment skills
    \item
      Database migration skills
    \end{itemize}
  \item
    Sharing skills across teams
  \item
    Hooks: automating actions around tool calls

    \begin{itemize}
    \tightlist
    \item
      Pre and post-execution hooks
    \item
      Notification hooks
    \item
      Validation and guard hooks
    \end{itemize}
  \end{itemize}
\item
  \texttt{RAG} (Retrieval-Augmented Generation) for Development

  \begin{itemize}
  \tightlist
  \item
    Why \texttt{RAG} matters for development

    \begin{itemize}
    \tightlist
    \item
      Overcoming context window limits
    \item
      Grounding \texttt{AI} in project-specific knowledge
    \item
      Reducing hallucinations
    \end{itemize}
  \item
    \texttt{RAG} architecture for codebases

    \begin{itemize}
    \tightlist
    \item
      Document chunking strategies for code
    \item
      Embedding models for source code
    \item
      Vector databases (Pinecone, \texttt{Weaviate}, \texttt{ChromaDB},
      \texttt{Qdrant})
    \item
      Hybrid search: vector + keyword
    \end{itemize}
  \item
    Indexing codebases

    \begin{itemize}
    \tightlist
    \item
      \texttt{File}-level vs function-level vs chunk-level indexing
    \item
      Handling multiple languages and file types
    \item
      Incremental indexing and keeping indexes fresh
    \item
      Metadata extraction and enrichment
    \end{itemize}
  \item
    \texttt{RAG} for documentation

    \begin{itemize}
    \tightlist
    \item
      Indexing internal wikis, READMEs, design docs
    \item
      Indexing \texttt{API} documentation
    \item
      Indexing Slack/Teams conversations and decision logs
    \end{itemize}
  \item
    Building a \texttt{RAG} pipeline

    \begin{itemize}
    \tightlist
    \item
      Ingestion pipeline design
    \item
      Query construction and rewriting
    \item
      Retrieval and re-ranking
    \item
      Context \texttt{assembly} and prompt construction
    \end{itemize}
  \item
    \texttt{RAG} evaluation and optimization

    \begin{itemize}
    \tightlist
    \item
      Measuring retrieval quality
    \item
      Measuring answer quality
    \item
      Iterating on chunking and embedding strategies
    \end{itemize}
  \end{itemize}
\item
  Prompt Engineering for Advanced Workflows

  \begin{itemize}
  \tightlist
  \item
    System prompts and persona design
  \item
    Chain-of-thought and step-by-step reasoning
  \item
    Few-shot and many-shot prompting
  \item
    Structured output: \texttt{JSON} mode, schemas, and constrained
    generation
  \item
    Prompt chaining and multi-stage pipelines
  \item
    Prompt templating and reuse
  \item
    Context management strategies

    \begin{itemize}
    \tightlist
    \item
      Priority-based context \texttt{assembly}
    \item
      Summarization and compression
    \item
      Sliding window approaches
    \end{itemize}
  \item
    Evaluating and iterating on prompts
  \item
    Prompt security: injection attacks and defenses
  \end{itemize}
\item
  Building Custom Agents

  \begin{itemize}
  \tightlist
  \item
    Agent \texttt{design\ patterns}

    \begin{itemize}
    \tightlist
    \item
      Single-turn vs multi-turn agents
    \item
      Supervisor and sub-agent architectures
    \item
      Parallel agent execution
    \item
      Human-in-the-loop patterns
    \end{itemize}
  \item
    Agent frameworks

    \begin{itemize}
    \tightlist
    \item
      \texttt{Anthropic\ Agent\ SDK}
    \item
      \texttt{LangChain} / \texttt{LangGraph}
    \item
      \texttt{CrewAI}
    \item
      \texttt{AutoGen}
    \item
      \texttt{OpenAI\ Agents\ SDK}
    \end{itemize}
  \item
    Implementing development agents

    \begin{itemize}
    \tightlist
    \item
      Code review agents
    \item
      Test generation agents
    \item
      Documentation agents
    \item
      Migration agents
    \item
      Monitoring and alerting agents
    \end{itemize}
  \item
    Agent memory and state

    \begin{itemize}
    \tightlist
    \item
      Short-term conversation memory
    \item
      Long-term persistent memory
    \item
      Project knowledge bases
    \item
      Learning from corrections
    \end{itemize}
  \item
    Agent evaluation and testing

    \begin{itemize}
    \tightlist
    \item
      Benchmarking agent performance
    \item
      Regression testing for agents
    \item
      Monitoring agent behavior in production
    \end{itemize}
  \end{itemize}
\item
  \texttt{AI} in the Development Lifecycle

  \begin{itemize}
  \tightlist
  \item
    \texttt{AI}-assisted planning and design

    \begin{itemize}
    \tightlist
    \item
      Requirements analysis
    \item
      Architecture exploration
    \item
      \texttt{API} design
    \item
      Data modeling
    \end{itemize}
  \item
    \texttt{AI}-powered code review

    \begin{itemize}
    \tightlist
    \item
      Automated review agents
    \item
      Style and convention enforcement
    \item
      Security scanning
    \item
      Performance analysis
    \end{itemize}
  \item
    \texttt{AI}-assisted testing

    \begin{itemize}
    \tightlist
    \item
      Unit test generation
    \item
      Integration test design
    \item
      Fuzzing and property-based testing
    \item
      Visual regression testing
    \end{itemize}
  \item
    \texttt{AI} in \texttt{CI/CD} pipelines

    \begin{itemize}
    \tightlist
    \item
      Automated PR checks with \texttt{AI}
    \item
      Intelligent test selection
    \item
      Deployment risk analysis
    \item
      Automated rollback decisions
    \end{itemize}
  \item
    \texttt{AI} for incident response

    \begin{itemize}
    \tightlist
    \item
      Log analysis and anomaly detection
    \item
      Root cause analysis
    \item
      Remediation suggestion
    \item
      Post-mortem generation
    \end{itemize}
  \end{itemize}
\item
  Advanced Topics

  \begin{itemize}
  \tightlist
  \item
    Multi-model strategies

    \begin{itemize}
    \tightlist
    \item
      Using different models for different tasks
    \item
      Model routing and fallback
    \item
      Cost optimization across models
    \end{itemize}
  \item
    Fine-tuning for development tasks

    \begin{itemize}
    \tightlist
    \item
      \texttt{When} to fine-tune vs prompt engineer
    \item
      Creating training datasets from codebases
    \item
      Evaluating fine-tuned models
    \end{itemize}
  \item
    Local models for development

    \begin{itemize}
    \tightlist
    \item
      Running models locally with Ollama, \texttt{LM\ Studio}, llama.cpp
    \item
      Privacy-sensitive development workflows
    \item
      Offline development with \texttt{AI}
    \item
      Hardware requirements and performance
    \end{itemize}
  \item
    \texttt{AI} governance and compliance

    \begin{itemize}
    \tightlist
    \item
      Code ownership and \texttt{AI}-generated code
    \item
      License implications
    \item
      Audit trails and reproducibility
    \item
      Organizational policies
    \end{itemize}
  \item
    Scaling \texttt{AI} development practices

    \begin{itemize}
    \tightlist
    \item
      Team onboarding and training
    \item
      Standardizing tool configurations
    \item
      Measuring and improving \texttt{AI} adoption
    \item
      Building internal \texttt{AI} development platforms
    \end{itemize}
  \end{itemize}
\end{itemize}

\subsection{Copyright}\label{copyright}

Mark Veltzer
\href{mailto:mark.veltzer@gmail.com}{\nolinkurl{mark.veltzer@gmail.com}},
© 2026

\end{document}
